% ============================================
% 四、程序完善题
% ============================================
\section{程序完善题}

% --- Q1: 字符类型判断 ---
\begin{frame}[fragile]{程序完善 第1题}{字符类型判断}
    \begin{exampleblock}{题目}
        输入一个字符，判断该字符是数字、字母、空格还是其他字符。
    \end{exampleblock}
    \begin{lstlisting}[language=C]
char ch;
ch = 【1】;                                    // getchar()
if (【2】)                                     // 判断字母
    printf("It is an English character\n");
else if (【3】)                                // 判断数字
    printf("It is a digit character\n");
else if (ch == ' ')
    printf("It is a space character\n");
else
    printf("It is other character\n");
\end{lstlisting}
    \pause
    \begin{alertblock}{答案}
        \begin{itemize}
            \item 【1】\texttt{getchar()} —— 从标准输入读取一个字符
            \item 【2】\texttt{(ch>='a'\&\&ch<='z')||(ch>='A'\&\&ch<='Z')}
            \item 【3】\texttt{(ch>='0')\&\&(ch<='9')}
        \end{itemize}
    \end{alertblock}
\end{frame}

% --- Q2: 查找数组元素 ---
\begin{frame}[fragile]{程序完善 第2题}{数组查找}
    \begin{exampleblock}{题目}
        在a数组中查找与x值相同的元素所在位置。
    \end{exampleblock}
    \begin{lstlisting}[language=C]
int a[10], i, x;
printf("input 10 integers:");
for (i = 0; i < 10; i++) scanf("%d", 【4】);  // &a[i]
printf("input the number you want to find x:");
scanf("%d", &x);
for (i = 0; i < 10; i++)
    if (【5】) break;                          // a[i]==x
if (i < 10)
    printf("the pos of x is: %d\n", i);
else
    printf("can not find x!\n");
\end{lstlisting}
    \pause
    \begin{alertblock}{答案}
        \begin{itemize}
            \item 【4】\texttt{\&a[i]} —— scanf需要变量\underline{地址}
            \item 【5】\texttt{a[i]==x} 或 \texttt{x==a[i]}
        \end{itemize}
    \end{alertblock}
\end{frame}

% --- Q3: 统计非负数 ---
\begin{frame}[fragile]{程序完善 第3题}{统计非负数}
    \begin{exampleblock}{题目}
        读入20个整数，统计非负数个数，并计算非负数之和。
    \end{exampleblock}
    \begin{lstlisting}[language=C]
int i, a[20], s, count;
s = 【6】;                              // count=0
for (i = 0; i < 20; i++) scanf("%d", &a[i]);
for (i = 0; i < 20; i++) {
    if (【7】) continue;                // a[i] < 0
    【8】                                // s += a[i];
    【9】                                // count++;
}
printf("s = %d\t count = %d\n", 【10】); // s, count
\end{lstlisting}
    \pause
    \begin{alertblock}{答案}
        \begin{itemize}
            \item 【6】\texttt{count=0}（注意s未初始化，但只累加非负数）
            \item 【7】\texttt{a[i]<0} —— 跳过负数
            \item 【8】\texttt{s+=a[i];} 或 \texttt{s=s+a[i];}
            \item 【9】\texttt{count++;} 或 \texttt{++count;} 或 \texttt{count+=1;}
            \item 【10】\texttt{s, count}（8、9空顺序可互换）
        \end{itemize}
    \end{alertblock}
\end{frame}
